\section{n1. Qualified framework, canonical approximation, and the exact smooth source}
\textit{Support: ENTIRE qualified G331 only. Reconstructed foundations here are not claimed as new results.}

Use the whole G331 grant, including its mandatory annotation, source ledger,
model scope and radial distinction. On the compact full section
$X=V_{\mathfrak a}(R)$ of every excellent reduced local germ, $C_R$ is the
continuous uniform completion of independently geometric evaluations:
actual negative effective finite REAL Cartier presentations on directed
ACTUAL primary blowups, with degree at least zero on EVERY contracted
proper integral curve before completion. Its proved comparison with the
negative primary/unit ideal cone, and the actual decreasing-$C$-net class
$D_R$, are used. All prime supports and quotient fraction fields are retained.
We use full finite-injection positive $C/D$ descent, completion and the
primary cocycle. No arbitrary-model cofinality or general envelope theorem
is part of G331. Bottom, zero/unit, reduction, and the empty dimension-zero
boundary retain their exact conventions.

All the following nonempty-section formulas assume positive-dimensional reduction;
the empty boundary has its unique empty function and envelope.
For finite continuous signed $f$, set
\[
 B=\max(0,-\min_X f),\quad
 A_m=\{b\in R:v(b)+mf(v)\ge0\text{ for EVERY }v\in X\},\quad
 A_0=R,\quad h_m=-v(A_m)/m.
\]
The inequalities, ultrametricity, and nonnegative values on $R$ show these
are ideals, contain $\mathfrak a^{\lceil mB\rceil}$, and satisfy
$A_mA_n\subset A_{m+n}$. If $f$ is negative somewhere they are proper
primary; otherwise they are the unit ideal. They are integrally closed:
if $b$ is integral over $A_m$ and $v(b)<-mf(v)$ is finite, in its monic
ideal-integrality equation the leading term has strictly smaller value than
all other terms. This is impossible. Infinite values already satisfy the
test. Every $h_m$ is finite continuous and below $f$. In the primary case,
$\operatorname{Bl}_{A_m}R$ has actual effective Cartier $E_m$ and
$\mathcal O(-E_m)$ relatively ample. Thus $-E_m/m$ has positive degree on
EVERY contracted integral curve before any limit. The unit case is zero.

Let $p_R(f)$ be the pointwise raw supremum of $C$ minorants. Rational
primary/unit profiles are uniformly dense in $C$ by G331. Approximate a
$C$ minorant, then subtract a small negative rational boundary constant.
A common denominator turns this into a primary ideal product. For a
rational competitor $-v(I)/r\le f$, $I\subset A_r$. This proves
\[
 p_R(f)=\sup_m h_m=\sup\{g:g\in D_R,\ g\le f\}.
\]
For the last equality, write a $D$ minorant as $\inf_i g_i$. For every
$\epsilon>0$ the directed open sets $\{g_i<f+\epsilon\}$ cover $X$.
Compactness selects one index throughout $X$, and $g_i-\epsilon$ is a
$C$ minorant. Let $\epsilon$ decrease. This includes partial bottom.
The constant $-B$ gives $-B\le p_R(f)\le\min(f,0)$; positivity or USC
of this raw supremum is not presumed.

Write $M_1=\|v(A_1)\|_\infty$, $H_N=\max_{1\le r\le N}h_r$.
If $m=kr+s$, $0\le s<r$, the grading gives
\[
 h_m\ge (kr/m)h_r-sM_1/m\ge h_r-rM_1/m,
 \qquad H_N-NM_1/m\le h_m\le p_R(f).
\]
Consequently every degree converges pointwise to the finite LSC function
$p_R(f)$. If $p_R(f)$ is continuous, the increasing open strict-test sets
$\{H_N>p_R(f)-\epsilon\}$ cover the compact section; one $N$ works
uniformly, then the displayed bound gives uniform convergence of $h_m$.
Conversely a uniform limit of $h_m$ is continuous and belongs to $C$.
A greatest $D$ minorant equals $p_R(f)$, so its USC, together with the
LSC already proved, makes it continuous. Hence greatest $C$ existence,
greatest $D$ existence, continuity of $p_R(f)$, and uniform canonical
convergence are equivalent. Shifting competitors by a negative constant
in both directions proves $\|p_R(f)-p_R(F)\|_\infty\le\|f-F\|_\infty$.

Here is the precise smooth source theorem used below, reconstructed from
complete immediate BFJ units freshly read in this run. Its domain is ONLY
a smooth affine CLOSED germ over an algebraically closed characteristic-zero
field, and its completion via G331. For dimension at least two, signed
primary model functions are determined on a fixed simple dual complex
(BFJ 1.14); they are uniformly dense (1.15). For such an obstacle $F$,
take the componentwise greatest weak-nef Weil divisor $W\le Z(F)$.
It exists: the negative constant barrier bounds a max-directed family from
below, $Z(F)$ bounds it from above, and its finite coefficient limit lies
in the closed nef cone. BFJ 3.11 provides substantive continuity on EACH
CLOSED simplex, not merely interior convex continuity. The extension
\[
 H(v)=\inf_{\pi} H_W(r_\pi v)
\]
uses ALL retractions. On any refinement of the fixed complex, divisorial
inequalities extend by closed-simplex continuity to every face, so
$-B\le H\le F$ on ALL semivaluations. Every $D$ competitor has its
nef coefficients below $W$, and the same continuity and all-retraction
formula give $g\le H$. This uses G331's exact full BFJ comparison;
identically bottom is handled separately.

The canonical completed primary/unit ideal $J_m$ is exactly
$L^\infty(mW)=L^\infty(mZ(F))$. Equality of the two ideals is BFJ 3.17:
$Z(b)/m$ is nef for a nonzero formal function $b$. Divisorial tests for
$Z(F)$ are equivalent to the full tests: truncate $(b)+\widehat{\mathfrak
m}^{N}$ with $N\ge mB$. Its primary/unit profile satisfies the tests on
divisorial points, hence everywhere by G331's all-support density and
continuity. Its value is at most $v(b)$, proving the assertion even when
$b$ is killed at a nonminimal support. G331 contracts these completed
primary ideals exactly to the germ.

For $W\ne0$, BFJ 3.14 gives uniform strict negativity, and thinness is
bounded on the FIXED complex. Thus $A_{\rm thin}\le-CW$ there for some
$C>0$. The WHOLE proof of BFJ 3.18 applies to real parameters: its (ii)
follows from $Z(b)\le tW+A_{\rm thin}$ on this complex, and its (iii)
follows by applying (i) to $tW$. Put $t=m+C$, rather than rounding away
this parameter, to obtain
\[
 (1+C/m)W\le W_m\le W,\quad W_m=m^{-1}Z(J_m),
 \qquad 0\le H-h_m\le CB/m
\]
on EVERY support. The extension from coefficients first uses closed-simplex
continuity, then the infimum over ALL retractions. In particular $H$ is a
uniform limit of actual ample primary profiles and belongs to $C$.
For $W=0$ the envelope is zero and $J_m$ is the unit ideal. Uniform
signed-model approximation and the proved Lipschitz bound now give a
continuous greatest $C/D$ minorant for EVERY finite continuous signed
obstacle on the exact smooth source. DVR and empty cases are direct.
General primary gauges transport this theorem with bounded positive
continuous denominators. Their model rate remains $K/m$ for a suitable
finite $K$; no rate for arbitrary continuous obstacles is asserted.

For another primary $\mathfrak b$, put $d(v)=v(\mathfrak b)>0$,
$T(v)=v/d(v)$ and $f_{\mathfrak b}(Tv)=f_{\mathfrak a}(v)/d(v)$.
All defining canonical inequalities divide by $d(v)$, so the ideals are
unchanged, and profiles and raw suprema divide by $d(v)$. Full completion
transport and exact primary contraction are those of G331, without a
claim about equivalence of arbitrary model categories.

The BFJ application uses the six complete bound carriers containing its
setting, simple modifications, 1.10--1.15, 2.1--2.8, 3.1--3.14 and
3.15--3.25, including complete proofs of 3.9--3.10 and 3.18.
Historical smooth resolution, multiplier subadditivity and smooth Izumi
remain disclosed published dependencies, not recursively fresh proofs.
No singular positive-characteristic transfer, normal-surface theorem,
CLD identification, products or energy is imported here.

For a reduced excellent curve, normalization has finitely many local DVR
centres. Each normalized centred DVR boundary is one point, since a
nonzero element is a uniformizer power times a unit. A regular conductor
element has finite order at every centre, so G331 makes restriction a
bijection of these full finite boundaries. Primary powers and units give
all nonpositive real values at each DVR point, and semilocal ideal gluing
gives their finite products. Full POSITIVE descent therefore identifies
$C_R$ with all finite nonpositive branch functions, $D_R$ with all extended
nonpositive ones, and the finite continuous obstacle envelope with
$\min(f,0)$ coordinatewise. Literal classical/CLD logarithmic radial
extensions remain distinct, including the qualified mixed-zero obstruction.
This is a comparison consequence, not a new result or radial identification.

Artifacts: notes/n1.tex, P12.3-quotients-and-wild-conductors.tex,
source-application-reading-ledger.md. Rerun: python build\_delivery.py;
python verify\_delivery.py. Artifact success does not certify this proof.
